\begin{abstract}
We formalize Sections 2.1--2.3 of Hassin and Peleg,
\emph{Distributed Probabilistic Polling and Applications to Proportionate Agreement},
Information and Computation 171 (2001), 248--268.
The synchronous process uses independent neighbor samples with directional weights.
On a finite connected nonbipartite undirected graph it reaches consensus with
probability one, and each color wins with its initial stationary vertex weight.
No path-space measure is needed: eventual probabilities are suprema of increasing
finite-time consensus probabilities.
\end{abstract}

\section{Shared foundation (generic library results)}
\begin{definition}[Finite weighted expectations]\label{def:weights}
\lean{Dynamics.Distribution,Dynamics.Kernel}\leanok
A distribution has nonnegative real weights summing to one. A kernel is a
family of such distributions, one for each state. Expectations are finite sums.
These are shared library definitions, not numbered results of the paper.
\end{definition}
\begin{theorem}[Stationary distribution exists]\label{thm:stationary}
\lean{Dynamics.Kernel.exists_stationary}\leanok\uses{def:weights}
Every finite nonempty stochastic matrix has a stationary distribution.
\end{theorem}
\begin{proof}\leanok
The shared library uses Ces\`aro averages of evolved distributions and
compactness of the finite probability simplex. This supplies the existence
fact recalled before Lemma 2.2; uniqueness is not required.
\end{proof}
\begin{theorem}[Finite-chain absorption]\label{thm:generic-absorption}
\lean{Dynamics.Kernel.exists_uniform_block,Dynamics.Kernel.geometric_blocks,Dynamics.Kernel.finite_absorption}\leanok\uses{def:weights}
An absorbing target accessible from every state of a finite chain admits a
uniform positive success probability in a fixed number of rounds. The survival
probability decays geometrically in blocks and tends to zero.
\end{theorem}
\begin{proof}\leanok
Take a common block length exceeding all finitely many access times, then a
maximum survival probability strictly below one. Iterate the resulting contraction.
\end{proof}

\section{Section 2.1: model}
\begin{definition}[Synchronous copying]\label{def:step}
\lean{Voter.step,Voter.round,Voter.transition}\leanok\uses{def:weights}
A configuration is $s:V\to C$. Independently sample $r(i)$ according to row
$H_i$, then set the next configuration to $i\mapsto s(r(i))$.
\end{definition}
\begin{lemma}[Configuration transition formula]\label{lem:product}
\lean{Voter.transition_product,Voter.transition_product_bool}\leanok\uses{def:step}
For configurations $s,t$, the transition probability is
\[P_{st}=\prod_i\sum_{j:s(j)=t(i)}H_{ij}.\]
For Boolean colors a white target vertex contributes $\sum_j H_{ij}s(j)$;
a black target contributes its complement, giving the paper's product formula.
\end{lemma}
\begin{proof}\leanok
The configuration distribution is the pushforward of the independently sampled
neighbors. Factor its equality indicator into a product over vertices.
\end{proof}

\section{Section 2.2: absorption and consensus probability}
\begin{lemma}[Monochromatic-edge propagation]\label{lem:propagation}
\lean{Voter.monochromatic_edge,Voter.propagate_region,Voter.possible_consensus}\leanok\uses{def:step}
On a connected nonbipartite graph every Boolean configuration has a possible
finite path to a constant configuration.
\end{lemma}
\begin{proof}\leanok
A coloring without a monochromatic edge would be a proper two-coloring.
Starting with a monochromatic edge, each vertex of its region can keep that
color by sampling an internal neighbor. Connectedness supplies an edge leaving
every proper region. One more vertex can copy across that edge. Induction on
the number of vertices outside the region gives consensus.
\end{proof}
\begin{lemma}[Lemma 2.1: absorption]\label{lem:21}
\lean{Voter.transition_constant,Voter.possible_positive,Voter.consensus_tendsto}\leanok\uses{lem:propagation,thm:generic-absorption,lem:product}
Assume $H_{ij}>0$ whenever $i,j$ are adjacent; extra support, such as
self-loops, is allowed. Both constant configurations
are absorbing, and the probability of remaining nonconstant tends to zero.
This is the finite-time survival formulation of the paper's transience claim.
\end{lemma}
\begin{proof}\leanok
Every possible copying round has positive probability. Thus graph propagation
gives accessibility. Apply the shared finite-chain absorption theorem.
\end{proof}
\begin{lemma}[Lemma 2.3: invariant stationary weight]\label{lem:23}
\lean{Voter.round_mass,Voter.iterate_mass}\leanok\uses{def:step,thm:stationary}
For any stationary distribution $\pi$, the expected white mass at every finite
time is the initial mass $\sum_i\pi_i s(i)$. No graph connectivity assumption
is imposed on this invariant.
\end{lemma}
\begin{proof}\leanok
A vertex's expected next color indicator is $\sum_jH_{ij}s(j)$.
Interchanging finite sums and using $\pi H=\pi$ proves one-step preservation;
induction proves the iterated statement.
\end{proof}
\begin{lemma}[Lemma 2.2: limiting identification]\label{lem:22}
\lean{Voter.whiteMass_bounds,Voter.whiteProbability_error,Voter.whiteProbability_tendsto}\leanok\uses{lem:21,lem:23}
The discrepancy between expected stationary white mass and all-white
probability lies between zero and nonconsensus probability. Consequently
all-white probability converges to initial stationary white mass.
The implementation combines the paper's limiting identification with the
invariant from Lemma 2.3 when stating the limit.
\end{lemma}
\begin{proof}\leanok
Pointwise, the white mass agrees with the all-white indicator on constant
configurations and lies in $[0,1]$ elsewhere. Take finite-time expectations
and squeeze the discrepancy using Lemma 2.1.
\end{proof}
\begin{theorem}[Theorem 2.1]\label{thm:21}
\lean{Voter.eventualColor,Voter.colorProbability_mono,Voter.consensus_probability}\leanok\uses{lem:22}
Define eventual all-white probability as the supremum of the increasing
finite-time all-white probabilities. It equals the initial stationary white mass.
\end{theorem}
\begin{proof}\leanok
Consensus is absorbing, so the probabilities increase and are bounded by one.
Their supremum is their limit. Identify this limit using Lemma 2.2.
\end{proof}
\begin{theorem}[Corollary 2.2 and regular graphs]\label{cor:22}
\lean{Voter.degree_stationary,Voter.uniform_consensus_probability,Voter.regular_consensus_probability}\leanok\uses{thm:21}
For uniform neighbor sampling, $\pi_i=d_i/(2|E|)$, so all-white consensus has
probability $\sum_{i:s(i)=1}d_i/(2|E|)$. On a regular graph this is the initial
fraction of white vertices.
\end{theorem}
\begin{proof}\leanok
The degree-sum formula normalizes these weights. In the stationarity equation,
the degree cancels the reciprocal sampling degree; undirected adjacency gives
the remaining column sum. Equal degrees then give uniform vertex weights.
\end{proof}

\section{Section 2.3: multiple colors}
\begin{theorem}[Consensus in a specified color]\label{thm:colors}
\lean{Voter.step_project,Voter.iterate_project,Voter.colorProbability_project,Voter.color_consensus_probability,Voter.sum_color_consensus_probability}\leanok\uses{thm:21}
For a finite color set and any specified color $c$, its eventual consensus
probability is $\sum_{i:s(i)=c}\pi_i$. These probabilities sum to one.
\end{theorem}
\begin{proof}\leanok
The color-indicator projection commutes with every copying step and with
iterated expectations. Its all-white event is exactly all-$c$ consensus.
Apply Theorem 2.1 to this Boolean projection.
\end{proof}

\section{Examples and scope}
\begin{lemma}[Triangle and a nonuniform matrix]\label{lem:examples}
\lean{Voter.Examples.triangle_one_white,Voter.Examples.biased_stationary}\leanok\uses{cor:22,lem:23}
One white vertex on a triangle wins with probability $1/3$.
The stochastic matrix whose two rows are $(1/3,2/3)$ has stationary
weights $(1/3,2/3)$, providing a nonuniform invariant example.
\end{lemma}
\begin{lemma}[Wright--Fisher sampling]\label{lem:wf}
\lean{Voter.Examples.wrightFisher,Voter.Examples.wrightFisher_stationary,Voter.Examples.wrightFisher_one_third}\leanok\uses{thm:21}
On three individuals, let every offspring pick a uniform parent, itself included.
This kernel has self-loops, so it is supported beyond the edges of the triangle.
Theorem~2.1 still applies, and one allele copy fixes with probability $1/3$.
\end{lemma}
\begin{lemma}[Bipartite two-cycle]\label{lem:cycle}
\lean{Voter.Examples.twoVertex_unique_neighbor,Voter.Examples.twoVertex_flip,Voter.Examples.twoVertex_cycle,Voter.Examples.twoVertex_never_consensus}\leanok\uses{def:step}
On the two-vertex graph, every vertex must copy the other vertex.
The alternating coloring flips each round and returns after two rounds;
neither coloring is constant. Thus connectedness alone is insufficient.
\end{lemma}
Convergence-time bounds, dynamic networks, and extremal coalition results
are future work beyond this milestone. Chernoff bounds remain in the
3-majority package until another theorem requires their extraction.
